\subsection{Acquisition costs and claim finance follow different clocks}
\label{sec:insurance-paa-cost-finance}

The one-year coverage example establishes PAA eligibility without shortening
the claim's settlement period. Acquisition costs add another choice. Under
the historical IFRS~17 text, costs incorporated into the PAA liability reduce
its opening balance and are subsequently amortised as an expense. Immediate
expense is permitted under 59(a) when each contract's coverage period at
initial recognition is no more than one year
\citep[IFRS~17.53--59]{eu20231803}.

Consider a separate eligible one-year group with premium 120 and acquisition
cash paid 12 at recognition. Opening cash and equity are each 200. Assume
uniform service and risk release, no investment component, tax or FX, and
the permitted LRC financing exception because the service-to-premium-due
interval is at most a year. No claims are incurred before coverage ends,
and the base case has no onerous indication. Cash after receipt and costs
is 308. Deferring the acquisition cost gives opening LRC \(120-12=108\)
and equity 200. Expensing it gives LRC 120 and equity 188.

At mid-year, both cases recognise revenue 60. In the deferred case, debit
acquisition expense 6 and credit LRC 6 for amortisation, then debit LRC 60
and credit revenue 60. LRC is \(108+6-60=54\), and cumulative profit
is 54. In the expensed case, LRC is 60 and cumulative profit is
\(60-12=48\). After the second half's service, both LRC balances are
zero and both cases have recognised revenue 120 and total acquisition
expense 12. Equal final costs do not justify erasing the interim choice.
Time allocation also requires the supplied uniform risk pattern; B126--127
instead uses expected incurred-service-expense timing when the risk-release
pattern differs significantly \citep[IFRS~17.55,59(a),B126--127]{eu20231803}.

An adverse mid-year branch requires a fresh onerous calculation. If current
FCF for the remaining coverage is 70 while deferred-cost LRC is 54, recognise
loss 16 and increase LRC to 70. Adding the entire FCF 70 to the existing
54 would overstate the liability by 54. The measurement must use the
discount treatment required by paragraph~57, including its connection to
the LIC simplification \citep[IFRS~17.57--58]{eu20231803}.

Return to the base case. At coverage end a claim is incurred, with expected
payment 48.4 two years later, current discount rate 10\% and separately
supplied risk adjustment 2. Initial LIC is \(48.4/1.1^2+2=42\);
cash 308 less LIC 42 leaves equity 266. The claim-incurrence rate of 10\%
differs from the group's assumed inception rate of 5\%. If the portfolio
elects to split LIC finance between profit and OCI, B133 uses the
claim-incurrence rate specified in B72(e)(iii), not the inception rate
\citep[IFRS~17.59(b),B72(e)(iii),B133]{eu20231803}.

One year after incurrence, suppose the current rate falls to zero and stays
there. The claim's current present value becomes 48.4, while a separate
schedule using the claim-incurrence rate reaches 44. Finance expense is
8.4 in total: 4 in profit and 4.4 in OCI. Assume the risk adjustment
releases 1 in each of the two years, wholly in the service result under
paragraph~81. First-year LIC is then 49.4 and equity 258.6. In the second
year, current claim value remains 48.4, so total finance is zero; the
claim-incurrence schedule reaches 48.4 through profit expense 4.4 and OCI
income 4.4. Release the remaining risk adjustment and pay the claim.
Final cash and equity are 259.6, equal to opening 200 plus premium 120
less acquisition costs 12 and settlement 48.4. Cumulative OCI is zero.
Using the inception rate would give first-year profit finance of only 2
and the wrong allocation, even if the total liability were correct.
